\scanpage{13}

\deleted{étant loc. connexe, et} ses composantes connexes \unsure{sont} $\in F_{\mathcal P}$ --- or $X_i$ étant un élément $\ne\varnothing$ minimal de $F_{\mathcal P}$, il s'ensuit que $X_i$ est connexe. Inversement, $3)\Longrightarrow d)$.

Soit en effet $A\in F_{\mathcal P}$ [loc fermé] \inserted{\unsure{pas néc}}. Quitte à remplacer $X$ par $A$, et la partition de $X$ par la partition induite [qui satisfait encore 1), 2), 3)] --- \deleted{\unsure{ceci}} c'est trivial sauf pour $2^\circ$, \inserted{mais} \deleted{\unsure{c'est vrai encore car $A$ est ouvert, mais}} si $X_i,X_j\subset A$, $i\ne j$, et $\overline X_i\deleted{\unsure{\cap}}\cap A=\overline X_j\cap A$, \deleted{comme \unsure{on a} $X_i$}\inserted{comme $X_i\subset\overline X_i\cap A$ et $X_j\subset\overline X_j\cap A$}, on aura $X_i\subset\overline X_j$ et $X_j\subset\overline X_i$ donc $\overline X_i=\overline X_j$ donc $X_i=\uline{X_j}$)

On est ramené à prouver \deleted{que $X$ est loc. connexe et} que les composantes connexes \deleted{\unsure{ouvertes}} de $X$ sont \deleted{\unsure{ouvertes et}} réunions de feuilles. Or les $X_i$ étant connexes, \deleted{les $\overline X_i$ aussi. Donc} \deleted{\unsure{et réunion des} $X_j$} les composantes connexes \inserted{de $X$} sont \inserted{bien} des réunions d'\unsure{ens} $\unsure{\overset{\deleted{-}}{X_i}|X_i}$, \deleted{\unsure{donc} (grâce à $2^\circ$) \unsure{disjoints} et \unsure{réunion des} \illegible{2}} Donc il n'y en a qu'un nb fini, \deleted{donc les composantes connexes sont \deleted{\unsure{f}} ouvertes, constituent une partition finie de $X$ \deleted{par des} par des fermés.}

Soit $Y=X/R\ (\simeq I)$ l'espace quotient. C'est un espace \uline{top} \uline{fini}. \deleted{\unsure{On a} $\unsure{\overline{\{i\}}}$} Si $i\in I$, $\unsure{\{i\}|\overline{\{i\}}}$ le plus petit

\mspage{14}
